\documentclass[11pt,a4paper]{article}
\usepackage[margin=1in]{geometry}
\usepackage{amsmath,amssymb,booktabs,longtable,hyperref}
\title{Human Compression: the audit of a fixed corpus}
\author{Human Compression project}
\date{11 October 2026}
\begin{document}
\maketitle

\section{The cutoff and the objects being compared}

New manuscript intake is closed. The baseline is the Git commit
\begin{center}\texttt{3b067cc2eba4ac2f2f91248448727639d8ac7c5a},\end{center}
also named \texttt{intake-2026-10-11}. Its catalogue contains 161 works, of which
102 have retrieved full texts, 61 have specified encodings, 48 have a
rederived result of some stated scope, and 27 have previously admitted
recoveries of specified principal results. These overlapping counts do
not describe 161 completed readings, 27 replacements of every assertion,
or 27 fully formalized papers.

The outstanding F015 review and L028 manuscript may finish because their
proof work was already underway. The bounded investigation of finite
minimal covers may record proved conditional criteria and unresolved
obstructions. N007 is not a further application in this intake. None of
these pending items is included in the 27 admitted principal contracts.
The exact cutoff records are in \texttt{audit/intake.json} and
\texttt{audit/source\_registry.json}.

For each source, the audit first fixes an edition and a list of results
to be recovered, with all parameter ranges and exceptional cases. A
correction to a false source statement is recorded as a correction, not
as a proof of the original statement. Recovery of the main results does
not by itself recover every auxiliary assertion, table or illustration.
The source and replacement must have the same declared mathematical
outputs before their lengths are compared.

\section{Theory, applications and proof closure}

A topic $A$ consists of a complete shared theory $U_A$ and complete
applications $A_1,\ldots,A_k$, corresponding to declared source contracts
$LA_1,\ldots,LA_k$. Each application must check the actual hypotheses of
the shared theorems and supply its local constructions, converse
arguments, exceptions, geometric identifications and any necessary
exhaustive computations. Similar wording is not a proof of common
structure. A paragraph saying that a theorem applies is sufficient only
when all its hypotheses and the claimed output have actually been
identified.

Existing proofs supply material for new common arguments. The objective is
to reprove the source results as consequences of a small theory, not to
shorten their wording. Moving a proof to another file earns no mathematical
credit. The earlier combined manuscripts are a proof bank, not the intended
replacement. The current position is as follows.

\begin{longtable}{p{0.12\textwidth}p{0.79\textwidth}}
\toprule Topic & Existing material and next obligation\\\midrule
A & A complete shared theory and the L032 application; the remaining
two-group results require separate complete applications.\\
B & A reviewed combined proof on forests and primitive coset geometries;
the shared theory and application papers still need separation.\\
C & A reviewed collection of halving and Petrial arguments; retain full
geometric identifications and computational classifications.\\
D & Reviewed marked-action and enumeration arguments; retain field and
group-specific classifications. L073 is partial, and L104 has a restricted
principal contract.\\
E & A new common theorem for string markings with proper exceptional
ends derives the intersection property instead of assuming it. Three
applications recover the two polytope classifications and the corrected
PSL maniplex results. Prior editions remain separate baselines.\\
F & The complete odd-kernel application L014; no separate complete shared
theory has yet been admitted.\\
G & Common flag-relation theory and the complete corrected L009
application; the retained Hermitian and classification proofs remain
dependencies.\\
Affine & Common affine theory and the L012 application, with the existing
Topic A theory and cubic arguments still required.\\\bottomrule
\end{longtable}

The first publication migration supplies ten editable TeX proof texts:
two in A, three in E, two in G, and three separate cover results. The
newly compiled PDFs have the same extracted text and page counts as the
reviewed PDFs. This checks textual preservation during conversion; it
does not establish new mathematics or complete the corpus audit.
The remaining proof texts and all required dependencies must undergo
the same treatment before a topic is declared ready for measurement.

\section{Attribution and formal verification}

Each application will have a machine-readable ledger connecting a source
result and its exact locator to the replacement theorem, its proof,
retained inputs, required programs, and any Lean declarations. The TeX
edition will distinguish borrowed results, contributions of this project,
recovery of several previous results, source corrections, computational
checks, formalized lemmas, and answers to previously posed questions.
An assertion of historical novelty requires a separate literature audit.

For every formal claim, the audit inspects the actual Lean statement,
including its hypotheses. A general transport lemma with an unproved
geometric hypothesis does not formalize the geometric theorem. The
ledger must name the written obligations left outside Lean. Finite
computations count as formal evidence only to the extent that their
certificates and the implication to the claimed theorem are checked.

The canonical library now contains 97 modules. The evidence consists of
a fresh run of 96 unchanged modules and a separate fresh compilation of
the new generic extension module, with nine checked declarations. This is
not a claim that all 97 were rebuilt in one run. All belong in the default
build. Source hashes, toolchain, dependency revisions and axiom reports
are recorded. One legacy
module has four native-evaluation axioms; that trust boundary remains
explicit. Existing limited formalizations will be extended where they
verify substantive steps in our own proofs. Neither successful compilation
nor the number of Lean declarations measures the fraction of a paper
that has been formalized.

\section{Length and reuse}

Fix the source contracts, representations, document boundaries and length
convention $c$ before reporting a factor. The intended principal comparison
uses the full source manuscripts once their declared main-result contracts
have been recovered. Ancillary material need not all be reproduced, but
the omitted results and any corrections must remain explicit. A comparison
restricted to matched proof passages is a separate sensitivity measure;
it must not silently replace the full-manuscript numerator.
For a complete topic, write
\[
 R_A=\frac{\sum_{i=1}^k c(LA_i)}
 {c(U_A)+\sum_{i=1}^k c(A_i)+c(P_A)},
 \qquad
 r_{A,i}\mid(U_A,P_A^{\mathrm{shared}})
 =\frac{c(LA_i)}{c(A_i)+c(P_{A,i})}.
\]
Here $P_A$ comprises required proof artifacts not already included in
the written costs, and $P_{A,i}$ is its application-specific part.
The conditional factor assumes that both the umbrella and shared required
artifacts have already been acquired. Shared artifacts are charged once
within the topic. Written-only and
artifact-inclusive comparisons are reported separately, using compatible
units rather than silently adding pages to code bytes. Classical inputs
and any unavailable historical artifact costs remain visible.

The complete cost of $U_B$ is charged in topic $B$, even if it overlaps
$U_A$. Conversely, $U_A$ is not charged again to every application in
topic $A$. Conditional application factors must display the cost of the
assumed umbrella. Missing costs are unknown, not zero. The previous
two-application Topic E measurement remains an edition-specific baseline:
$1.297\times$ for complete written manuscripts, or $1.131\times$ after
charging our complete required program. The latter is a known-cost
comparison because historical program costs are unknown. The earlier
$0.835\times$ used selected principal-proof passages. None of these
figures measures the best possible compression or the subsequent common
proof edition. Its three-source comparison uses the same three contracts
on both the old and new replacement sides.

A separate sequential account follows the reuse discipline employed in
the GKM project. Let $K_{i-1}$ be the already paid-for library and $D_i$
the full dependency closure of the next application, in an additive
representation of identified proof units. Then
\[
 \Delta_i=c(D_i\setminus K_{i-1}),\qquad
 K_i=K_{i-1}\cup D_i.
\]
Hypothesis checks, new interfaces and required programs belong to $D_i$.
The fixed graph and an explicit admissible order determine this account.
Without unrecorded revisions, the increments telescope to the cost of
the acquired union; no earlier theory disappears from the total.
This sequential account does not alter the independent topic costs.

Both the source and replacement require theorem-dependency graphs with
the same output contracts. Complete proof units, including exposition
and interfaces, receive costs. A citation alone does not establish that
a result is used; a cited theorem that is actually used remains an explicit
retained-input edge.
Publish per-paper costs, cumulative costs and sequential acquisition
costs. A sawtooth pattern must emerge from those measurements. It is
neither an assumption of the audit nor an established conclusion of the
current work.

\section{Dependency nodes and their costs}

For each matched result, construct a directed acyclic graph whose arrows
run from a prerequisite to the argument using it. Numbered theorems,
lemmas, propositions and corollaries give a first level of proof units.
Substantive unnumbered arguments receive explicit units too; definitions,
context, references and programs are distinguished from mathematical
proof nodes. A proof node contains its full statement, proof and local
hypothesis checks. Repeating a statement does not create a second result,
but the printed repetition still has a cost.

For an output $t$, let $D(t)$ contain $t$ and all its prerequisites. We
report the number $N(t)$ of mathematical proof nodes in $D(t)$, the list
of external inputs, and the known weighted cost
\[
 W_{\rm known}(t)=\sum_{v\in D(t),\,w(v)\text{ known}}w(v).
\]
The omitted unknown weights are listed explicitly; this sum is never
presented as an unconditional total. Statements retained from a classical
library have their printed statement cost and a separate unknown proof
cost. Required programs are charged with the proof they support.

Each measured document is partitioned into disjoint text spans. A span
has exactly one cost owner. Nodes can own several spans; node weights
are their combined costs. Context outside the logical closure remains
charged in the full-manuscript comparison. This makes the weighted graph
additive without replacing the primary full-document numerator by a
selected-proof numerator. Edges record actual use with a locator, not
merely the appearance of a citation.

Node counts depend on granularity. A theorem containing six cases must
not earn a sixfold saving simply because another author numbered the
cases as lemmas. We therefore retain the named-result count and expose
substantive branches where bundling matters. Parent and child descriptions
are two views of the same proof: costs belong to the leaves and are not
charged again at their parents. The weighted total must be invariant
under such regrouping. Counts, weighted costs and their granularity are
reported together; no raw node ratio is called a complexity ratio.

For several outputs the acquired closure is their union, so a common
proof is charged once within the topic. After acquiring $U_A$, the next
application pays its new closure and its own hypothesis checks. Source
papers receive the same reuse treatment for earlier results they actually
invoke. Separate topics still pay their whole umbrellas independently,
as in the preceding section. The paired graphs should identify where a
new argument eliminates a chain of old arguments, where it merely moves
one, and where an unchanged external theorem is still required.

\section{Questions and publication}

Answers to open questions are assessed separately from recovery of known
results. Each entry states the exact question and edition, the answer's
scope, proof dependencies, computational and formal evidence, and limits
of the priority search. They are not inserted into a historical-length
numerator. In particular an absent historical proof has no measured
length, so an infinite compression factor is not an empirical length
ratio. The project may identify a qualitative gain beyond recovery while
keeping this numerical distinction explicit.

TeX and Lean are the editable mathematical sources; PDFs are readers and
JSON supplies ledgers. Markdown stays in local working directories.
The 992 historical Markdown files remain bound to their exact Git objects,
with their paths and hashes in \path{audit/archive_markdown.json}.
The current verifier ignores local notes and checks the archived bytes
separately from the current tracked publication. The old replay programs
run from a complete checkout of the pinned intake commit. They are not
represented as fresh replays of the reorganized publication.

The next audit completes the matched Topic E dependency closure, then
reworks the remaining topic papers from existing proofs. A topic closes
only after independent mathematical review, complete source attribution,
the declared Lean and computation checks, and a measured denominator
containing every required component. This cutoff changes the work from
acquisition to proof consolidation and assessment; it makes no new claim
of recovered results, solved questions or compression.
\end{document}
